\entry{Semana del 15/09/2025.} 

\section{Lunes 15/09/2025 - Charla en la RAFA -SUF} 
Hoy fue la Charla de la RAFA. Hubieron dos preguntas: 

\begin{itemize}
	\item Cómo se compara la escala de las mediciones del sensor en la cuba chica con las de una boya en el océano. 
	\item Si la capilaridad se debe a las fuerzas intermoleculares del agua o a qué. 
\end{itemize}



\section{Viernes 19/09/2025} 
El Miércoles de la semana pasada terminaron de imprimir en el Taller el nuevo modelo de carcasa del sensor capacitivo, que permite conectarlo directamente a un perfil de 40x40 mm. Juan y Mateo pasaron a buscarlo el Martes. La impresión parece haber salido correctamente, solo que vino con todo el soporte y algunos filamentos colgando. Sí parecía que se corrió o movió durante las primeras capas ya que una de las ``alas" para atornillar ambas mitades estaba un poco deformada, pero los tornillos pasaron igualmente. 

Lo monté, dejando más largo para que sea más fácil sumergirlo en la cuba grande desde los laterales y verifiqué que funcionase bien la coneción al perfil.  

\begin{figure}[th!]
	\centering
	\includegraphics[width=0.567\linewidth]{"Figures/15_09_2025/Nuevo modelo del sensor montado"}
	\caption{Nuevo odelo de carcasa del sensor capacitivo ya montado en perfil de 40x40 mm. Solo cambió el conector respecto al modeo anterior, que tenía dos entradas de tornillos para que conectase al tornillo micrométrico.} 
	\label{fig:nuevo-modelo-del-sensor-montado}
\end{figure}


Faltaría recubrirlo con el esmalte aislante. 